\documentclass[lang=cn,newtx,10pt,scheme=chinese]{elegantbook}
\usepackage{paracol} % 或 parallelx

\title{希亚瓦萨号列车}
\subtitle{Pullman Car Hiawatha}

\author{桑顿·怀尔德}
\institute{抱了又抱戏剧社}

\setcounter{tocdepth}{3}

\cover{cover.jpg}

% 本文档命令
\usepackage{array}
\newcommand{\ccr}[1]{\makecell{{\color{#1}\rule{1cm}{1cm}}}}

\newcommand{\pair}[2]{%
\switchcolumn* #2\par
\switchcolumn  #1\par
}

% 修改标题页的橙色带
\definecolor{customcolor}{RGB}{32,178,170}
\colorlet{coverlinecolor}{customcolor}
\usepackage{cprotect}

\begin{document}

\maketitle
\frontmatter

\tableofcontents

\mainmatter

\part{剧本}


\begin{paracol}{2}


\pair{
    {\huge\textbf{人物}}
}{
    {\huge\textbf{CHARACTERS}}
    \newline
    \newline
}

\pair{舞台监督}{THE STAGE MANAGER}

\pair{
    \textit{三号包厢}：
    \begin{itemize}[leftmargin=5em, labelsep=0.6em, itemsep=0.2em]
        \item 一个疯女人，丘吉尔夫人
        \item 男性乘客，摩根先生
        \item 女性乘客，一名有经验的护士
    \end{itemize}
}{
    \textit{Compartment Three}:
    \begin{itemize}[leftmargin=5em, labelsep=0.6em, itemsep=0.2em]
        \item AN INSANE WOMAN, Mrs. Churchill
        \item A MALE ATTENDANT, Mr. Morgan
        \item THE FEMALE ATTENDANT, A trained nurse
    \end{itemize}
}

\pair{\textit{二号包厢}：菲利普·米尔布里}{\textit{Compartment Two}: PHILIP}

\pair{\textit{一号包厢}：哈瑞特·米尔布里，菲利普的年轻太太}{\textit{Compartment One}: HARRIET, Philip's young wife}

\pair{\textit{一号下铺}：一个少女}{\textit{Lower One}: A MAIDEN LADY}

\pair{\textit{三号下铺}：一名中年医生}{\textit{Lower Three}: A MIDDLE-AGED DOCTOR}

\pair{\textit{五号下铺}：一个肥胖却可亲的女人}{\textit{Lower Five}: A STOUT, AMIABLE WOMAN OF FIFTY}

\pair{\textit{七号下铺}：一名工程师，比尔，准备前往加利福尼亚}{\textit{Lower Seven}: AN ENGINEER, Bill, going to California}

\pair{\textit{九号下铺}：一名工程师，弗雷德}{\textit{Lower Nine}: AN ENGINEER, Fred}

\pair{乘务员：哈里森}{THE PORTER, Harrison}


\pair{俄亥俄州格洛佛角：形象是一个满脸笑容的男孩}{GROVER'S CORNERS, OHIO}
\pair{田野：形象是一个穿衬衫的人}{THE FIELD}
\pair{一个流浪汉}{THE TRAMP}
\pair{俄亥俄州帕克斯堡：形象是一个农民的妻子和三个年轻人}{PARKERSBURG, OHIO}
\pair{一个工人：克鲁格先生，鬼魂}{THE WORKMAN, Mr. Krüger, a ghost}
\pair{另一个工人：看守}{THE WORKER, a watchman}
\pair{天气：形象是一名机械工程师}{A MECHANIC}
\pair{时间：十点，十一点，十二点(形象是三个漂亮女人)}{The Hours: TEN O'CLOCK, ELEVEN O'CLOCK, TWELVE O'CLOCK}
\pair{星球：土星，金星，木星，地球}{The Planets: SATURN, VENUS, JUPITER, EARTH}
\pair{
    大天使：加百列，米迦勒
}{
    The Archangels: GABRIEL, MICHAEL
    \newline
    \newline
    \newline
}

\pair{
    {\huge \textbf{场景}}
}{
    {\huge \textbf{SETTING}}
    \newline
    \newline
    \newline
}

\pair{
    一辆从纽约开往芝加哥的卧铺列车，1930年12月。
}{
    A Pullman car making its way from New York to Chicago, December 1930.
}

\pair{
    \textit{舞台后方是一处类似阳台或吊桥、逃生梯的地方，向两侧延展至观众视线之外。两段楼梯自上方延伸到舞台，没有其他额外的布景。}
}{
    \textit{At the back of the stage is a balcony or bridge or runway leading out of sight in both directions. Two flights of stairs descend from it to the stage. There is no further scenery.}
}

\pair{
    \textit{幕启，舞台监督正用粉笔在脚灯附近的地板上画线。}
}{
    \textit{At the rise of the curtain The Stage Manager is making lines with a piece of chalk on the floor of the stage by the footlights.}
}

\pair{
    \textbf{舞台监督}：这是一辆卧铺列车的布局，这趟车名叫希亚瓦萨号，12月21日，正从纽约开往芝加哥。你们的左手边是三个包厢，这边是过道和五个下铺床位。卧铺已经满员，上铺下铺都有人，不过在这出戏里，我们只聚焦于稍远处这一侧的下铺。现在是九点半，几个铺位都已经收拾妥当，大部分乘客都躺在绿色帘子后面的小床上。他们正在脱了鞋往地上扔，要么是在费劲地脱着裤子，或是正在琢磨着，把贵重物品藏在枕头里晚上会不会掉出来。好啦！大家都上来吧！
}{
    THE STAGE MANAGER: This is the plan of a Pullman car. Its name is Hiawatha and on December twenty-first it is on its way from New York to Chicago. Here at your left are three compartments. Here is the aisle and five lowers. The berths are all full, uppers and lowers, but for the purposes of this play we are limiting our interest to the people in the lower berths on the further side only. The berths are already made-up. It is half past nine. Most of the passengers are in bed behind the green curtains. They are dropping their shoes on the floor, or wrestling with their trousers, or wondering whether they dare hide their valuables in the pillow slips during the night. All right! Come on, everybody!
    \newline
    \newline
}

\pair{
    \textit{演员们搬着椅子上场。他们将两张椅子对着摆放在粉笔划定的区域，代表自己的铺位。他们侧身坐着，将脚搭在另一把椅子上，躺下就可以睡觉。包厢内的乘客姿势也是一样。从左到右我们依次看到：三号包厢，二号包厢，一号包厢，一号下铺，三号下铺，五号下铺，七号下铺，九号下铺。}
}{
    \textit{The actors enter carrying chairs. Each improvises his berth by placing two chairs “facing one another” in his chalk-marked space. They then sit in one chair, profile to the audience, and rest their feet on the other. This must do for lying in bed. The passengers in the compartments do the same.}
    \newline
    \newline
    \newline
}

\pair{
    \textbf{一号下铺}：乘务员，记得六点一刻叫我起来。
}{
    LOWER ONE: Porter, be sure and wake me up at quarter of six.
}

\pair{
    \textbf{乘务员}：好的，女士。
}{
    THE PORTER: Yes, ma'am.
}

\pair{
    \textbf{一号下铺}：我肯定也睡不着，不过六点一刻还是叫我一下。
}{
    LOWER ONE: I know I shan't sleep a wink, but I want to be told when it's quarter of six.
}

\pair{
    \textbf{乘务员}：好的，女士。
}{
    THE PORTER: Yes, ma'am.
}

\pair{
    \textbf{七号下铺}（\textit{从帘子后面伸出头}）：哎！乘务员！哎！你这个灯到底怎么打开啊？
}{
    LOWER SEVEN (Putting his head through the curtains): Hsst! Porter! Hsst! How the hell do you turn on this other light?
}

\pair{
    \textbf{乘务员}（\textit{有些烦躁}）：估计坏了，唉。你只能用另外一头的。
}{
    THE PORTER (\textit{Fussing with it}): I'm afraid it's outta order, suh. You'll have to use the other end.
}

\pair{
    \textbf{舞台监督}（\textit{用假嗓子扮演一位上铺的女士}）：我想问一下，车上有没有好心人能给我一点阿司匹林？
}{
    THE STAGE MANAGER (\textit{Falsetto, substituting for some woman in an upper berth}): May I ask if someone in this car will be kind enough to lend me some aspirin?
}

\pair{
    \textbf{乘务员}（\textit{忙碌地来回奔走}）：好的，女士。
}{
    THE PORTER (\textit{Rushing about}): Yes, ma'am.
}

\pair{
    \textbf{九号下铺}（\textit{其中一名工程师从过道下来直接就去了五号下铺}）：对不起，女士，对不起。走错了。
}{
    LOWER NINE (\textit{One of the engineers, descending the aisle and falling into Lower Five}): Sorry, lady, sorry. Made a mistake.
}

\pair{
    \textbf{五号下铺}（\textit{嘟囔}）：这辈子都没见过这种事儿！
}{
    LOWER FIVE (\textit{Grumbling}): Never in all my born days!
}

\pair{
    \textbf{一号下铺}（\textit{小声地}）：乘务员！乘务员！
}{
    LOWER ONE (\textit{In a shrill whisper}): Porter! Porter!
}

\pair{
    \textbf{乘务员}：你好，女士。
}{
    THE PORTER: Yes, ma'am.
}

\pair{
    \textbf{一号下铺}：我的热水袋在漏水，你得把它弄走。今晚上我没有热水袋用了，真惨。
}{
    LOWER ONE: My hot water bag's leaking. I guess you'll have to take it away. I'll have to do without it tonight. How awful!
}

\pair{
    \textbf{五号下铺}（\textit{对上铺的乘客厉声说}）：小伙子，你别那么多事行吗？不然我就要去列车长那儿投诉了。
}{
    LOWER FIVE (\textit{Sharply to the passenger above her}): Young man, you mind your own business, or I'll report you to the conductor.
}

\pair{
    \textbf{舞台监督}（\textit{扮演五号上铺}）：对不起，女士。我真不是故意的，我的背带掉下去了，我就是想把它们拉上来。
}{
    THE STAGE MANAGER (\textit{Substituting for Upper Five}): Sorry, ma'am, I didn't mean to upset you. My suspenders fell down and I was trying to catch them.
}

\pair{
    \textbf{五号下铺}：来，给你。赶紧睡觉吧。感觉今晚上所有人都在往我这儿来。（\textit{把头伸出去}）乘务员！乘务员！发发善心帮我拿杯水行吗？我口干得厉害。
}{
    LOWER FIVE: Well, here they are. Now go to sleep. Everybody seems to be rushing into my berth tonight. (\textit{She puts her head out}) Porter! Porter! Be a good soul and bring me a glass of water, will you? I'm parched.
}

\pair{
    \textbf{九号下铺}：比尔！（\textit{没有人回应。}）比尔！
}{
    LOWER NINE: Bill! (\textit{No answer.}) Bill!
}

\pair{
    \textbf{七号下铺}：哈？你要啥？
}{
    LOWER SEVEN: Yea? Wha'd'ya want?
}

\pair{
    \textbf{九号下铺}：给我递本杂志行吗？
}{
    LOWER NINE: Slip me one of those magazines, willya?
}

\pair{
    \textbf{七号下铺}：你要哪本？
}{
    LOWER SEVEN: Which one d'ya want?
}

\pair{
    \textbf{九号下铺}：都行，《侦探故事集》吧。都可以。
}{
    LOWER NINE: Either one. Detective Stories. Either one.
}

\pair{
    \textbf{七号下铺}：弗雷德，我正看着《侦探故事集》呢，有个故事看到一半。
}{
    LOWER SEVEN: Aw, Fred. I'm just in the middle of one of 'em in Detective Stories.
}

\pair{
    \textbf{九号下铺}：行吧，那我看《西部冒险》——谢了。
}{
    LOWER NINE: That's all right. I'll take the Western.—Thanks.
}

\pair{
    \textbf{舞台监督}：（\textit{对演员们}）好了！嘘！嘘！嘘！（\textit{对观众}）现在，我要让大家听听他们在想什么。
}{
    THE STAGE MANAGER (\textit{To the actors}): All right! —Sh! Sh! Sh! (\textit{To the audience}) Now I want you to hear them thinking.
    \newline
    \newline
}

\pair{
    \textit{片刻停顿后，演员们纷纷开始低语，发出轻柔的沙沙声响。观众会逐一听到每个人的独白。}
}{
    \textit{There is a pause and then they all begin a murmuring-swishing noise, very soft. In turn each one of them can be heard above the others.}
    \newline
    \newline
}

\pair{
    \textbf{五号下铺}（\textit{五十岁的女人}）：我想想，给小宝贝的玩偶，给玛丽耶塔的鞋子，给赫伯特的钢笔。还有给乔治订阅《时代周刊》……
}{
    LOWER FIVE (\textit{The Woman of Fifty}): Let's see: I've got the doll for the baby. And the slip-on for Marietta. And the fountain pen for Herbert. And the subscription to Time for George...
}

\pair{
    \textbf{七号下铺}（\textit{比尔}）：我的天！莉莉安，如果你和我想象的不一样，我真不知道该怎么办了。告诉一个女人你专门跑到加利福尼亚去见她，大概不是一个明智的做法。我还是编个故事好了，就说我是去出差什么的。以前我有为了什么人这么激动这么纠结过吗？好吧，玛莎算一个，但那次不一样。我还是看看书吧，不然我该疯了。“你怎么知道访客离开的时候是十点钟？”侦探问道。“因为十点钟的时候，”女孩回答，“我会去关储藏室和走廊的灯。我下楼的时候听见先生在门口和什么人在聊天，我听见他说‘晚安……’”妈呀，我根本没在看书，我满脑子都是莉莉安。她的黄色头发，她的眼睛！
}{
    LOWER SEVEN (\textit{Bill}): God! Lillian, if you don't turn out to be what I think you are, I don't know what I'll do. —I guess it's bad politics to let a woman know that you're going all the way to California to see her. I'll think up a song-and-dance about a business trip or something. Was I ever as hot and bothered about anyone like this before? Well, there was Martha. But that was different. I'd better try and read or I'll go cuckoo. “How did you know it was ten o'clock when the visitor left the house?” asked the detective. “Because at ten o'clock,” answered the girl, “I always turn out the lights in the conservatory and in the back hall. As I was coming down the stairs I heard the master talking to someone at the front door. I heard him say, ‘Well, good night...’” —Gee, I don't feel like reading; I'll just think about Lillian. That yellow hair. Them eyes!...
}

\pair{
    \textbf{三号下铺}（\textit{医生正在给自己朗读着医学期刊上那些令人毛骨悚然的文章，时不时地停顿一下，发出一声充满质疑的“所以呢？”。}）
}{
    LOWER THREE (\textit{The Doctor reads aloud to himself the most hair-raising material from a medical journal, every now and then punctuating his reading with an interrogative “So?”})
}

\pair{
    \textbf{一号下铺}（\textit{少女}）：今天晚上我是睡不着了，最好现在就意识到这个问题。真不知道那个漏水的热水袋应该放到哪儿。我现在就靠着床右边躺好，深呼吸，想一些美好的事情，说不定我能打一会盹儿。
}{
    LOWER ONE (\textit{The Maiden Lady}): I know I'll be awake all night. I might just as well make up my mind to it now. I can't imagine what got hold of that hot water bag to leak on the train of all places. Well now, I'll lie on my right side and breathe deeply and think of beautiful things, and perhaps I can doze off a bit.
}

\pair{
    \textbf{九号下铺}（\textit{弗雷德}）：那是我干过的最疯狂的事情，直接让我的人生倒退三年。如果我就待在这里，可能已经攒了三万美金了。我有什么资格和苏联人签合同？见了鬼了，我还觉得挺有意思。有意思，真是见鬼！直接倒退了三年。我都不知道公司还会不会要我。我太嫩了，真的。我就是还没长大。
}{
    LOWER NINE (\textit{Fred}): That was the craziest thing I ever did. It's set me back three whole years. I could have saved up thirty thousand dollars by now, if I'd only stayed over here. What business had I got to fool with contracts with the goddam Soviets. Hell, I thought it would be interesting. Interesting, what the hell! It's set me back three whole years. I don't even know if the company'll take me back. I'm green, that's all. I just don't grow up.
    \newline
    \newline
}

\pair{
    \textit{舞台监督大步走向他们，一边举着手喊道“嘘”，大家停止了低语。}
}{
    \textit{The Stage Manager strides toward them with lifted hand, crying, “Hush,” and their whispering ceases.}
    \newline
    \newline
}

\pair{
    \textbf{舞台监督}：可以了——稍等，乘务员！
}{
    THE STAGE MANAGER: That'll do!—Just one minute. Porter!
}

\pair{
    \textbf{乘务员}（\textit{出现在左侧}）：来了先生！
}{
    THE PORTER (\textit{Appearing at the left}): Yessuh.
}

\pair{
    \textbf{舞台监督}：轮到你了。（\textit{乘务员很害羞}）。你不想表达吗？你有权利表达。
}{
    THE STAGE MANAGER: It's your turn to think. (\textit{The Porter is very embarrassed.}) Don't you want to? You have a right to.
}

\pair{
    \textbf{乘务员}（\textit{想要表达自己的想法，又很害羞，十分纠结}）：啊，啊，我就是在想我在芝加哥的家还有……我的保险。
}{
    THE PORTER (\textit{Torn between the desire to release his thoughts and his shyness}): Ah...ah...I'm only thinkin' about my home in Chicago and ... and my life insurance.
}

\pair{
    \textbf{舞台监督}：就是这样。
}{
    THE STAGE MANAGER: That's right.
}

\pair{
    \textbf{乘务员}：谢谢……谢谢您。
}{
    THE PORTER: ...Well, thank you... Thank you.
    \newline
    \newline
}

\pair{
    \textit{乘务员悄悄走开，突然的自我意识让他很激动，脸红得厉害。}
}{
    \textit{The Porter slips away, blushing violently, in an agony of self-consciousness and pleasure.}
    \newline
    \newline
}

\pair{
    \textbf{舞台监督}（\textit{对观众}）：这是个好小伙。哈里森是个好小伙子，就是害羞。（\textit{对演员们}）现在，该包厢的各位了。
}{
    THE STAGE MANAGER (To the audience): He's a good fellow, Harrison is. Just shy. (\textit{To the actors again}) Now the compartments, please.
    \newline
    \newline
}

\pair{
    \textit{卧铺区域光变暗。菲利普正站在他的包厢和妻子包厢的连接门处。}
}{
    \textit{The berths fall into shadow. Philip is standing at the door connecting his compartment with his wife's.}
    \newline
    \newline
}

\pair{
    \textbf{菲利普}：我的天使，你还好吗？
}{
    PHILIP: Are you all right, angel?
}

\pair{
    \textbf{哈瑞特}：还好。不知道晚饭的时候我到底怎么了。
}{
    HARRIET: Yes. I don't know what was the matter with me during dinner.
}

\pair{
    \textbf{菲利普}：需要我把门关上吗？
}{
    PHILIP: Shall I close the door?
}

\pair{
    \textbf{哈瑞特}：你看能不能用一把椅子顶着门，这样它能半开着，而且不会来回撞着响。
}{
    HARRIET: Do see whether you can't put a chair against it that will hold it half open without banging.
}

\pair{
    \textbf{菲利普}：好了。晚安，我的天使。你要是睡不着就叫我，我们可以起来打牌。
}{
    PHILIP: There.—Good night, angel. If you can't sleep, call me and we'll sit up and play Russian Bank.
}

\pair{
    \textbf{哈瑞特}：你肯定是想到了那段糟糕的日子，咱们有一个礼拜，每天晚上都那么坐着……不过我知道我今晚该好好睡。车轮子的声音听起来甚至有些悦耳，有回家的感觉。我们到哪个州了？
}{
    HARRIET: You're thinking of that awful time when we sat up every night for a week... But at least I know I shall sleep tonight. The noise of the wheels has become sort of nice and homely. What state are we in?
}

\pair{
    \textbf{菲利普}：我们正在穿过俄亥俄，马上就要到印第安纳了。
}{
    PHILIP: We're tearing through Ohio. We'll be in Indiana soon.
}

\pair{
    \textbf{哈瑞特}：我知道，那儿的小镇到处都有骑马墩。
}{
    HARRIET: I know those little towns full of horse blocks.
}

\pair{
    \textbf{菲利普}：我们一早就到芝加哥了。你踏踏实实睡觉，到时候我会叫你。
}{
    PHILIP: Well, we'll reach Chicago very early. I'll call you. Sleep tight.
}

\pair{
    \textbf{哈瑞特}：踏踏实实睡，亲爱的。
}{
    HARRIET: Sleep tight, darling.
    \newline
    \newline
}

\pair{
    \textit{菲利普回到自己的包厢。三号包厢里，男性乘客让自己的椅背斜靠着墙，正在抽烟。那位有经验的护士正在缝袜子，疯女人将额头顶在窗户上，正好直视着观众的方向。}
}{
    \textit{Philip returns to his own compartment. In Compartment Three, the male attendant tips his chair back against the wall and smokes a cigar. The trained nurse knits a stocking. The insane woman leans her forehead against the windowpane, that is, stares into the audience.}
    \newline
    \newline
}

\pair{
    \textbf{疯女人}（\textit{拖沓而抱怨的语气，但又充满不确定}）：别带我走。别带我走。
}{
    THE INSANE WOMAN (\textit{Her words have a dragging, complaining sound, but lack any conviction}): Don't take me there. Don't take me there.
}

\pair{
    \textbf{女性乘客}：亲爱的，你要不要躺一会儿？
}{
    THE FEMALE ATTENDANT: Wouldn't you like to lie down, dearie?
}

\pair{
    \textbf{疯女人}：我想下车，我想回纽约。
}{
    THE INSANE WOMAN: I want to get off the train. I want to go back to New York.
}

\pair{
    \textbf{女性乘客}：要不要我再给你梳梳头？多舒服啊。
}{
    THE FEMALE ATTENDANT: Wouldn't you like me to brush your hair again? It's such a nice feeling.
}

\pair{
    \textbf{疯女人}（\textit{走到门口}）：我要下车，我要开门。
}{
    THE INSANE WOMAN (\textit{Going to the door}): I want to get off the train. I want to open the door.
}

\pair{
    \textbf{女性乘客}（\textit{握住她的一只手}）：太大声了！你这样会把所有人都吵醒的。过来，我给你讲个故事，关于我们这趟旅程的目的地。
}{
    THE FEMALE ATTENDANT (\textit{Taking one of her hands}): Such a noise! You'll wake up all the nice people. Come and I'll tell you a story about the place we're going to.
}

\pair{
    \textbf{疯女人}：我不想到那儿去。
}{
    THE INSANE WOMAN: I don't want to go to that place.
}

\pair{
    \textbf{女性乘客}：那里很不错！到处都是草地和花园，我从来没见过那么好的地方，真的很好。
}{
    THE FEMALE ATTENDANT: Oh, it's lovely! There are lawns and gardens everywhere. I never saw such a lovely place. Just lovely.
}

\pair{
    \textbf{疯女人}（\textit{躺在床上}）：有玫瑰花吗？
}{
    THE INSANE WOMAN (\textit{Lies down on the bed}): Are there roses?
}

\pair{
    \textbf{女性乘客}：当然了！红的黄的白的，到处都有。
}{
    THE FEMALE ATTENDANT: Roses! Red, yellow, white... just everywhere.
}

\pair{
    \textbf{男性乘客}（\textit{短暂的停顿后}）：应该到克利夫兰了。
}{
    THE MALE ATTENDANT (\textit{After a pause}): That musta been Cleveland.
}

\pair{
    \textbf{女性乘客}：我之前有个病人在克利夫兰。糖尿病。
}{
    THE FEMALE ATTENDANT: I had a case in Cleveland once. Diabetes.
}

\pair{
    \textbf{男性乘客}（\textit{又是一个短暂的停顿}）：真希望我现在有一台收音机，听听广播对他们有好处。我以前有个病人，一秒钟都离不开收音机。
}{
    THE MALE ATTENDANT (\textit{After another pause}): I wisht I had a radio here. Radios are good for them. I had a patient once that had to have the radio going every minute.
}

\pair{
    \textbf{女性乘客}：收音机是好东西。我那个结了婚的侄子有一台，他也总是放广播，感觉很不错。
}{
    THE FEMALE ATTENDANT: Radios are lovely. My married niece has one. It's always going. It's wonderful.
}

\pair{
    \textbf{疯女人}（\textit{半起身}）：我长得不漂亮。我不像她那么漂亮。
}{
    THE INSANE WOMAN (\textit{Half rising}): I'm not beautiful. I'm not beautiful as she was.
}

\pair{
    \textbf{女性乘客}：噢，我觉得你漂亮！漂亮——摩根先生，你不觉得丘吉尔夫人很漂亮吗？
}{
    THE FEMALE ATTENDANT: Oh, I think you're beautiful! Beautiful.—Mr. Morgan, don't you think Mrs. Churchill is beautiful?
}

\pair{
    \textbf{男性乘客}：噢，很好看的！您像个电影明星，丘吉尔夫人。
}{
    THE MALE ATTENDANT: Oh, fine lookin'! Regular movie star, Mrs. Churchill.
    \newline
    \newline
}

\pair{
    \textit{疯女人疑惑地看看他们，又躺了下去。哈瑞特小声呻吟着，闷声咳了一下。她伸手摸索，按了响铃。乘务员来敲门。}
}{
    \textit{The Insane Woman looks inquiringly at them and subsides. Harriet groans slightly. Smothers a cough. She gropes about with her hand and finds the bell. The Porter knocks at her door.}
    \newline
    \newline
}

\pair{
    \textbf{哈瑞特}（\textit{低声细语}）：请进。请先把通往我先生房间的门关上。轻点，轻一点。
}{
    HARRIET (\textit{Whispering}): Come in. First, please close the door into my husband’s room. Softly. Softly.
}

\pair{
    \textbf{乘务员}（\textit{伤感地}）：好的，女士。
}{
    THE PORTER (\textit{A plaintive porter}): Yes, ma'am.
}

\pair{
    \textbf{哈瑞特}：乘务员，我不太舒服。我病了，需要看医生。
}{
    HARRIET: Porter, I'm not well. I'm sick. I must see a doctor.
}

\pair{
    \textbf{乘务员}：怎么了，女士，可这儿没有医生……
}{
    THE PORTER: Why ma'am, they ain't no doctor...
}

\pair{
    \textbf{哈瑞特}：有的，我刚才吃过晚餐出来的时候，看到那边有个男人，他在看医学杂志。你去叫他一下。
}{
    HARRIET: Yes, when I was coming out from dinner I saw a man in one of the seats on that side, reading medical papers. Go and wake him up.
}

\pair{
    \textbf{乘务员}（\textit{十分惊讶}）：女士，我不能去把别人叫醒。
}{
    THE PORTER (\textit{Flabbergasted}): Ma'am, I cain't wake anybody up.
}

\pair{
    \textbf{哈瑞特}：你可以，乘务员。乘务员先生，咱们别争了，我病得很厉害，是我的心脏。去把他叫起来，跟他说我心脏有问题。
}{
    HARRIET: Yes, you can. Porter. Porter. Now don't argue with me. I'm very sick. It's my heart. Wake him up. Tell him it's my heart.
}

\pair{
    \textbf{乘务员}：好的，女士。
}{
    THE PORTER: Yes, ma'am.
    \newline
    \newline
}

\pair{
    \textit{他走过过道，开始摇晃三号下铺的肩膀。}
}{
    \textit{He goes into the aisle and starts pulling the shoulder of the man in Lower Three.}
    \newline
    \newline
}

\pair{
    \textbf{三号下铺}：哎，哎，怎么了？到站了？（\textit{乘务员对他小声说话}）我马上过来——乘务员，是年轻人还是老人？
}{
    LOWER THREE: Hello. Hello. What is it? Are we there? (\textit{The Porter mumbles to him.}) I'll be right there. —Porter, is it a young woman or an old one?
}

\pair{
    \textbf{乘务员}：不知道，先生。我觉得像是有点老，嗯，也不算太老。
}{
    THE PORTER: I dunno, suh. I guess she's kinda old, suh, but not so very old.
}

\pair{
    \textbf{三号下铺}：你跟她说我马上就到，让她安静地躺着。
}{
    LOWER THREE: Tell her I'll be there in a minute and to lie quietly.
    \newline
    \newline
}

\pair{
    \textit{乘务员走进哈瑞特的包厢，她已经把头转了过去。}
}{
    \textit{The Porter enters Harriet's compartment. She has turned her head away.}
    \newline
    \newline
}

\pair{
    \textbf{乘务员}：女士，他马上就到。他让您安静地躺着。
}{
    THE PORTER: He'll be here in a minute, ma'am. He says you lie quiet.
    \newline
    \newline
}

\pair{
    \textit{三号下铺踉跄着走过过道，一边嘟囔着:“妈的，都是鞋！”}
}{
    \textit{Lower Three stumbles along the aisle muttering: “Damn these shoes!”}
    \newline
    \newline
}

\pair{
    \textbf{某人的声音}：这辆车能不能安静一会儿？
}{
    SOMEONE'S VOICE: Can't we have a little quiet in this car, please?
}

\pair{
    \textbf{九号下铺}（\textit{弗雷德}）：哦，闭上你的嘴！
}{
    LOWER NINE (\textit{Fred}): Oh, shut up!
    \newline
    \newline
}

\pair{
    （\textit{医生绕过乘务员，走进哈瑞特的包厢。他凑到她身边，整个人把她挡了起来}）
}{
    \textit{The Doctor passes The Porter and enters Harriet's compartment. He leans over her, concealing her by his stooping figure.}
    \newline
    \newline
}

\pair{
    \textbf{三号下铺}：乘务员，她死了。车上有人和她是一起的吗？
}{
    LOWER THREE: She's dead, Porter. Is there anyone on the train traveling with her?
}

\pair{
    \textbf{乘务员}：有的，先生。她丈夫就在这儿。
}{
    THE PORTER: Yessuh. Dat's her husband in dere.
}

\pair{
    \textbf{三号下铺}：蠢货！你刚才怎么不叫他？我进去跟他说吧。
}{
    LOWER THREE: Idiot! Why didn't you call him? I'll go in and speak to him.
    \newline
}

\pair{
    \textit{舞台监督走上前。}
}{
    \textit{The Stage Manager comes forward.}
    \newline
}

\pair{
    \textbf{舞台监督}：好了，关于车厢内就看这么多，目前来说已经够了。现在该从地理上、气候上、天文学上、神学上来了解下我们所处的位置。希亚瓦萨号列车，10点10分，1930年12月21日，好的。（\textit{后方高台上出现了一些身影}）不不，还没到星球出场的时候呢，时间也不该出场。
}{
    THE STAGE MANAGER: All right. So much for the inside of the car. That'll be enough of that for the present. Now for its position geographically, meteorologically, astronomically, theologically considered. Pullman Car Hiawatha, ten minutes of ten. December twenty-first, 1930. All ready. (\textit{Some figures begin to appear on the balcony.}) No, no. It's not time for The Planets yet. Nor The Hours.
    \newline
}

\pair{
    \textit{他们退场。}
}{
    \textit{They retire}
    \newline
}

\pair{
    \textit{舞台监督拍拍手，一个穿着工作裤、满脸笑容的小男孩从左侧卧铺后面上场。}
}{
    \textit{The Stage Manager claps his hands. A grinning boy in overalls enters from the left behind the berths.}
    \newline
}

\pair{
    \textbf{俄亥俄州格洛佛角}（\textit{口音傻愣愣的，就像是在学校周末活动中背诵文章}）：我表演的是俄亥俄州格洛佛角，这里有821人。“最坏的人也有很多好处，最好的人也有很多坏处，我们最好不要说人长短。”——罗伯特·路易斯·斯蒂文森。
}{
    GROVER'S CORNERS, OHIO (\textit{In a foolish voice as though he were reciting a piece at a Sunday school entertainment}): I represent Grover's Corners, Ohio. Eight hundred twenty-one souls. “There's so much good in the worst of us and so much bad in the best of us, that it ill behooves any of us to criticize the rest of us.” Robert Louis Stevenson. Thankya.
    \newline
    \newline
}

\pair{
    \textit{谢谢。他笑着从右边下场。从刚刚同一方向，一个穿衬衫的人上场，这是一片田野。}
}{
    \textit{He grins and goes out right. Enter from the same direction somebody in shirt sleeves. This is a field.}
    \newline
    \newline
}

\pair{
    \textbf{田野}：我表演的是田野，就是你们从俄亥俄格洛佛角到俄亥俄帕克斯堡经过的这片地方。在这片土地上，有51只金花鼠，206只田鼠，6条蛇，几百万只小虫子、昆虫、蚂蚁和蜘蛛。它们都在冬眠。“有什么能如六月的一天一样珍稀？如果有一天，完美的日子会到来。”——《郎法尔爵士的异象》，威廉·库伦——我想说的是詹姆斯·罗素·洛威尔。谢谢。
}{
    THE FIELD: I represent a field you are passing between Grover's Corners, Ohio, and Parkersburg, Ohio. In this field there are fifty-one gophers, two hundred and six field mice, six snakes and millions of bugs, insects, ants and spiders. All in their winter sleep. “What is so rare as a day in June? Then, if ever, come perfect days.” The Vision of Sir Launfal, William Cullen—I mean James Russell Lowell. Thank you.
    \newline
    \newline
}

\pair{
    \textit{田野下场，流浪汉上场。}
}{
    \textit{Exit. Enter a tramp.}
    \newline
    \newline
}

\pair{
    \textbf{流浪汉}：我想说一下，我是个流浪汉，一直跟着这趟车，希亚瓦萨号，所以我也应该出个场。我从纽约罗切斯特出发，要去伊利诺伊的朱丽叶特。很多很多人组成了这个世界。“去往曼德勒的路上，飞鱼嬉戏，黎明的曙光如雷鸣般，在海湾那边的中国升起。”——弗兰克·W·瑟维斯。感觉有点冷啊。谢谢。
}{
    THE TRAMP: I just want to tell you that I'm a tramp that's been traveling under this car, Hiawatha, so I have a right to be in this play. I'm going from Rochester, New York, to Joliet, Illinois. It takes a lotta people to make a world. “On the road to Mandalay, where the flying fishes play and the sun comes up like thunder, over China, ’cross the bay.” Frank W. Service. It's bitter cold. Thank you.
    \newline
    \newline
}

\pair{
    \textit{流浪汉下场，一个温柔的农妇带着三个纤瘦的年轻人上场。}
}{
    \textit{Exit. Enter a gentle old farmer's wife with three stringy young people.}
    \newline
    \newline
}

\pair{
    \textbf{俄亥俄州帕克斯堡}：我表演的是俄亥俄州帕克斯堡。这里有2604人。我曾亲眼见证酒精是如何带来毁灭的，希望在这个美丽的国家，不再有人沾染到哪怕一滴诅咒。（\textit{起了个节奏，大家都摇摆着唱了起来}）“扔出救生索！快扔出救生索！有人今日正在沉沦……”
}{
    PARKERSBURG, OHIO: I represent Parkersburg, Ohio. Twenty-six hundred and four souls. I have seen all the dreadful havoc that alcohol has done and I hope no one here will ever touch a drop of the curse of this beautiful country. (\textit{She beats a measure and they all sing unsteadily:}) “Throw out the lifeline! Throw out the lifeline! Someone is sinking today-ay...”
    \newline
    \newline
}

\pair{
    \textit{舞台监督顺势摆摆手，让他们下场。一个工人上场。}
}{
    \textit{The Stage Manager waves them away tactfully. Enter a workman.}
    \newline
    \newline
}

\pair{
    \textbf{工人}：Ich bin der Arbeiter der hier sein Leben verlor. Bei der Sprengung fur diese Brücke über die Sie in dem Moment fahren——（\textit{火车头发出呜呜声响，列车正在跨越栈桥}）——erschlug mich ein Felsblock. Ich spiele jetzt als Geist in diesem Stück mit. “Vor sieben und achtzig Jahren haben unsere Väter auf diesem Kontinent eine neue Nation hervorgebracht…”
}{
    THE WORKMAN: Ich bin der Arbeiter der hier sein Leben verlor. Bei der Sprengung für diese Brücke über die Sie in dem Moment fahren—(\textit{The engine whistles for a trestle crossing})—erschlug mich ein Felsbrock. Ich spiele jetzt als Geist in diesem Stück mit. “Vor sieben und achtzig Jahren haben unsere Väter auf diesem Kontinent eine neue Nation hervorgebracht...”
}

\pair{
    \textbf{舞台监督}（\textit{向观众解释}）:抱歉，他说的是德语。他说，他是一个鬼魂，死在修筑栈桥的过程中，就是希亚瓦萨号列车正经过的这座桥。（\textit{火车又发出呜呜声响}）。所以他想在这部剧中亮个相。他是在爆破的时候被一块大石头砸中的。他的座右铭大家应该听过：“六十七年前，我们的先辈们在这片大陆上创立了一个新国家……”就是这些。谢谢你，克鲁格先生。
}{
    THE STAGE MANAGER (\textit{Helpfully, to the audience}): I'm sorry; that's in German. He says that he's the ghost of a workman who was killed while they were building the trestle over which the car Hiawatha is now passing—(\textit{The engine whistles again})—and he wants to appear in this play. A chunk of rock hit him while they were dynamiting. —His motto you know: “Three score and seven years ago our fathers brought forth upon this continent a new nation dedicated...” and so on. Thank you, Mr. Krüger.
    \newline
    \newline
}

\pair{
    \textit{鬼魂下场，另一个工人上场。}
}{
    \textit{Exit the ghost. Enter another worker.}
    \newline
    \newline
}

\pair{
    \textbf{另一个工人}：我是俄亥俄州帕克斯堡附近一座塔楼的看守。我就是想说一下，我还醒着，看到这趟列车途中信号一切正常。希望大家旅途愉快。“假如你能保持冷静，即使众人都失去理智并且归咎于你……”——鲁德亚德·吉卜林。谢谢。
}{
    THE WORKER: I'm a watchman in a tower near Parkersburg, Ohio. I just want to tell you that I'm not asleep and that the signals are all right for this train. I hope you all have a fine trip. “If you can keep your head when all about you are losing theirs and blaming it on you...” Rudyard Kipling. Thank you.
    \newline
    \newline
}

\pair{
    \textit{下场。舞台监督走上前。}
}{
    \textit{He exits. The Stage Manager comes forward.}
    \newline
    \newline
}

\pair{
    \textbf{舞台监督}：好了，这部分也差不多了。该轮到天气了。
}{
    THE STAGE MANAGER: All right. That'll be enough of that. Now the weather.
}

\pair{
    \textit{一名机械工程师上场。}
}{
    \textit{Enter a mechanic.}
}

\pair{
    \textbf{机械工程师}：现在是11摄氏度，北风转西北风，风速57。一股低气压正由萨斯喀彻温向东海岸移动。明日天气寒冷，中西部各州和纽约北部会有降雪。
}{
    A MECHANIC: It is eleven degrees above zero. The wind is north-northwest, velocity: fifty-seven. There is a field of low barometric pressure moving eastward from Saskatchewan to the eastern coast. Tomorrow it will be cold with some snow in the middle western states and northern New York.
    \newline
}

\pair{
    \textit{下场。}
}{
    \textit{He exits}
    \newline
}

\pair{
    \textbf{舞台监督}：好的，该轮到时间们了。（\textit{向观众解释}）分钟是说闲话的人，小时是哲学家，年份都是神学家。小时中有一个例外——十二点钟也是一位神学家。准备好，十点钟！
}{
    THE STAGE MANAGER: All right. Now for The Hours. (\textit{Helpfully to the audience}) The minutes are gossips; the hours are philosophers; the years are theologians. The hours are philosophers with the exception of Twelve O'clock who is also a theologian. —Ready Ten O'clock!
    \newline
}

\pair{
    \textit{时间们的表演者都是美丽的女孩，打扮得好像以利·维达的画作《阿特拉斯的七个女儿》中一样。每个女孩都手拿金色的罗马数字，她们在后方的高台上，从右向左缓缓穿行而过。}
}{
    \textit{The Hours are beautiful girls dressed like Elihu Vedder's Pleiades. Each carries a great gold Roman numeral. They pass slowly across the balcony at the back, moving from right to left.}
    \newline
}

\pair{
    \textbf{舞台监督}：十点钟，你在干什么？亚里士多德？
}{
    THE STAGE MANAGER: What are you doing, Ten O'clock? Aristotle?
}

\pair{
    \textbf{十点}：不对，华士本先生，我是柏拉图。
}{
    TEN O'CLOCK: No, Plato, Mr. Washburn.
}

\pair{
    \textbf{舞台监督}：好——“您难道不相信，这样……”
}{
    THE STAGE MANAGER: Good.—"Are you not rather convinced that he who thus..."
}

\pair{
    \textbf{十点}：“难道不相信，他这样把美视作唯一可被看见的事物，是非常可贵的？因为他所看到的不是影子，而是真实……”
}{
    TEN O'CLOCK: "Are you not rather convinced that he who thus sees Beauty as only it can be seen will be specially favored? And since he is in contact not with images but with realities..."
    \newline
}

\pair{
    \textit{她继续低语着自己的篇章，十一点上场。}
}{
    \textit{She continues the passage in a murmur as Eleven O'clock appears}
    \newline
}

\pair{
    \textbf{十一点}：“我，爱比克泰德，一个糟老头子，除了赞美上帝，还能做些什么呢？如果我是一只夜莺，就会像夜莺一样歌唱。如果我是一只天鹅，就会像天鹅一样高歌。而我只是理性的造物……
}{
    ELEVEN O'CLOCK: "What else can I, Epictetus, do, a lame old man, but sing hymns to God? If then I were a nightingale, I would do the nightingale's part. If I were a swan, I would do a swan's. But now I am a rational creature..."
    \newline
}

\pair{
    \textit{她的声音渐小，成了低语。十二点上场。}
}{
    \textit{Her voice also subsides to a murmur. Twelve O'clock appears}
    \newline
}

\pair{
    \textbf{舞台监督}：好——十二点，你呢？
}{
    THE STAGE MANAGER: Good. —Twelve O'clock, what have you?
}

\pair{
    \textbf{十二点}：奥古斯丁和他的母亲。
}{
    TWELVE O'CLOCK: Saint Augustine and his mother.
}

\pair{
    \textbf{舞台监督}：那么——“我们说：如果肉身的躁动安静了……”
}{
    THE STAGE MANAGER: So. —"And we began to say: If to any the tumult of the flesh were hushed..."
}

\pair{
    \textbf{十二点}：“我们说：如果肉身的躁动安静了，地球、水、空气的样子也安静了……”
}{
    TWELVE O'CLOCK: "And we began to say: If to any the tumult of the flesh were hushed; hushed the images of earth; of waters and of air..."
}

\pair{
    \textbf{舞台监督}：快一点——“天国的地极也安静了。”
}{
    THE STAGE MANAGER: Faster. —"Hushed also the poles of Heaven."
}

\pair{
    \textbf{十二点}：“是的，灵魂也安静了。”
}{
    TWELVE O'CLOCK: "Yea, were the very soul to be hushed to herself."
}

\pair{
    \textbf{舞台监督}：大点声，福斯特小姐。
}{
    THE STAGE MANAGER: A little louder, Miss Foster.
}

\pair{
    \textbf{十二点}（\textit{声音大了一些}）：“梦与想象中的启示也会安静……”
}{
    TWELVE O'CLOCK (A little louder): "Hushed all dreams and imaginary revelations..."
}

\pair{
    \textbf{舞台监督}：（\textit{挥手示意她们下场}）好了，好了。该轮到星球了。1930年12月21日。
}{
    THE STAGE MANAGER (Waving them back): All right. All right. Now The Planets. December twenty-first, 1930, please.
    \newline
    \newline
}

\pair{
    \textit{时间们放松下来，退回到右侧的更衣室。星球们出现在高台上，有些人站在楼梯上。他们没有台词，但是会发出声音。有人是脉搏跳动的声音，有人是敲击的声音，有人用口哨从低到高再从高到低吹着音阶。土星的声音缓慢而单调，在两个低音上来回重复。}
}{
    \textit{The Hours unwind and return to their dressing rooms at the right. The Planets appear on the balcony. Some of them take their place halfway on the steps. These have no words, but each has a sound. One has a pulsating, zinging sound. Another has a thrum. One whistles ascending and descending scales. Saturn does a slow, obstinate humming sound on two repeated low notes:}
    \newline
    \newline
}

\pair{
    \textbf{舞台监督}：土星，大点声。金星，声音再高点。好的，该木星了。地球。（\textit{转身面向火车上的铺位}）来吧，各位。这是地球的声音。\textit{镇、工人，还有其他人都出现在舞台边缘。乘客们开始低语，发出他们“思考的声音”。}来，格洛佛角，帕克斯堡，你们都有份。看守，流浪汉，这是地球的声音。
}{
    THE STAGE MANAGER：Louder, Saturn. —Venus, higher. Good. Now, Jupiter. —Now the Earth. (\textit{The Stage Manager turns to the beds on the train.}) Come, everybody. This is the Earth's sound. (\textit{The towns, workmen, etc., appear at the edge of the stage. The passengers begin their “thinking” murmur.}) Come, Grover's Corners. Parkersburg. You're in this. Watchman. Tramp. This is the Earth’s sound.
    \newline
}

\pair{
    \textit{他像一个乐团的指挥一样指挥着大家。小镇和工人都在念着自己的座右铭。疯女人突然激动得哭了起来。她起身，向舞台监督挥着手。}
}{
    \textit{He conducts it as the director of an orchestra would. Each of the towns and workmen does his motto. The Insane Woman breaks into passionate weeping. She rises and stretches out her arms to The Stage Manager.}
    \newline
}

\pair{
    \textbf{疯女人}：使用我吧，让我做些什么。
}{
    THE INSANE WOMAN: Use me. Give me something to do.
    \newline
}

\pair{
    \textit{舞台监督快步走向她，对她低语几句，将她交给后面的监护人，她并没有得到安慰。}
}{
    \textit{He goes to her quickly, whispers something in her ear, and leads her back to her guardians. She is unconsoled.}
    \newline
}

\pair{
    \textbf{舞台监督}：嘘！嘘嘘！大天使来了。（\textit{对观众}）现在我们可以从神学视角来看看希亚瓦萨号列车了。
}{
    THE STAGE MANAGER: Now shh—shh—shh! Enter The Archangels. (\textit{To the audience}) We have now reached the theological position of Pullman Car Hiawatha.
    \newline
}

\pair{
    \textit{小镇和工人们都下场了。星球们也已经退场，仍然能听到微弱的音乐声。两个穿着蓝色套装的年轻人从高台上出现，沿着右边的楼梯走下来。他们经过乘客身边，乘客们正在自己的梦中低语。
    \newline
    加百列让米迦勒看看比尔，比尔正挑着眉露出笑容。他们在五号下铺停住了脚步，米迦勒发出“嗯——嗯”式的赞同声音。
    \newline
    角色们在睡梦中的言论并非全部清晰易懂，会被他们的叹息、呻吟或低语淹没，但我们隐约能听到一些。}
}{
    \textit{The towns and workmen have disappeared. The Planets, offstage, continue a faint music. Two young men in blue serge suits enter along the balcony and descend the stairs at the right. As they pass each bed the passenger talks in his sleep.
    \newline
    Gabriel points out Bill to Michael who smiles with raised eyebrows. They pause before Lower Five, and Michael makes the sound of assent that can only be rendered “Hn-Hn.”
    \newline
    The remarks that the characters make in their sleep are not all intelligible, being lost in the sound of sigh or groan or whisper by which they are conveyed. But we seem to hear:}
}

\pair{
    \textbf{九号下铺}（\textit{大声地}）：有些人比另一些人动作更慢，仅此而已。
}{
    LOWER NINE (\textit{Loud}): Some people are slower than others, that's all.
}

\pair{
    \textbf{七号下铺}（\textit{比尔}）：你知道，这不是玩笑，我会努力的。
}{
    LOWER SEVEN (\textit{Bill}): It's no fun, y'know. I'll try.
}

\pair{
    \textbf{五号下铺}（\textit{买圣诞礼物的女士，语速很快}）：你们肯定知道，我都给你们准备好了，年年都是如此。
}{
    LOWER FIVE (\textit{The lady of the Christmas presents, rapidly}): You know best, of course. I'm ready whenever you are. One year's like another.
}

\pair{
    \textbf{一号下铺}：我可以教别人缝纫，我会缝东西。
}{
    LOWER ONE: I can teach sewing. I can sew.
    \newline
}

\pair{
    \textit{他们来到哈瑞特的包厢，疯女人坐了起来，对他们讲话。}
}{
    \textit{They approach Harriet’s compartment. The Insane Woman sits up and speaks to them.}
    \newline
}

\pair{
    \textbf{疯女人}：是我吗？（\textit{大天使们摇摇头。}）我在这里干等着有什么用处呢？我对一切心怀感恩，我对自己变得比从前更好而感恩。那些过去的事，那些可怕的事，已经不再纠缠我了。我的思想中，那些重担已经卸下了——但是不再有人能理解我了。我完全能够理解我自己，但是我说什么别人都听不懂——我还要等待？
}{
    THE INSANE WOMAN: Me? (\textit{The Archangels shake their heads.}) What possible use can there be in my simply waiting? —Well, I'm grateful for anything. I'm grateful for being so much better than I was. The old story, the terrible story, doesn't haunt me as it used to. A great load seems to have been taken off my mind. —But no one understands me any more. At last I understand myself perfectly, but no one else understands a thing I say. — So I must wait?
    \newline
}

\pair{
    \textit{大天使们微笑着点头。他们无可奈何地微笑，表明这是一个决定。}
}{
    \textit{The Archangels nod, smiling. Resignedly, and with a smile that implies their complicity}
    \newline
}

\pair{
    \textbf{疯女人}（\textit{继续说着}）：好吧，你们来决定，我会尽我的力。只是人们都太幼稚、太荒唐了。他们毫无逻辑，很疯狂，像小孩儿一样，他们从没经受过苦难。
}{
    THE INSANE WOMAN (\textit{continuing}): Well, you know best. I'll do whatever is best; but everyone is so childish, so absurd. They have no logic. These people are all so mad... These people are like children; they have never suffered.
    \newline
}

\pair{
    \textit{她重回睡梦之中。大天使们站在哈瑞特身旁。医生刚把菲利普带到另一个包厢，正在认真地跟他小声讲话。哈瑞特面朝着墙壁，她微微转过头，面对着车顶开始说话。}
}{
    \textit{She returns to her bed and sleeps. The Archangels stand beside Harriet. The Doctor has drawn Philip into the next compartment and is talking to him in earnest whispers. Harriet’s face has been toward the wall; she turns it slightly and speaks toward the ceiling.}
    \newline
}

\pair{
    \textbf{哈瑞特}：我在那儿不会快乐的，让我就这样死在这里吧。我属于这里，我就待在我的小房子旁边，待在菲利普身边就很好。你知道的，我在那儿不会快乐的。（\textit{加百列靠近她对她说话，一个短暂的停顿，她突然激动得哭了起来。}）我没有脸面去那儿，我什么都不曾做过，我这一生全都虚度了。我脾气暴躁，性格阴郁，什么都不懂，我没资格踏进那里一步。（\textit{大天使们再次对她低语。}）怎么可能就这样宽恕我呢，我不该这么轻易就得到宽恕。我希望得到惩罚，除非我得到足够的惩罚，我哪里都不去。我想要彻底地解脱——在受到惩罚之后，我想做一个全新的人。（\textit{大天使们对她低声说话时，她的双脚慢慢踩在地上。}）没有人可以为我受罚，我想自己去面对，我不想任何人为我受罚……（\textit{他们对她说话，她坐在那里盯着自己的鞋，努力地思考着。}）这不公平，我应该自己去受苦——如果我可以承受那么重的痛苦。（\textit{她微笑了起来}）噢，我真的不配！你们知道的，我只是个愚蠢的人，我只是一个普通的人——可现在这么美好的事情要来临了。你们真的要让我去？你们真的想要我去？（\textit{他们牵着她走过列车的过道。}）
}{
    HARRIET: I wouldn't be happy there. Let me stay dead down here. I belong here. I shall be perfectly happy to roam about my house and be near Philip. —You know I wouldn't be happy there. (\textit{Gabriel leans over and whispers into her ear. After a short pause she bursts into fierce tears.}) I'm ashamed to come with you. I haven't done anything. I haven't done anything with my life. Worse than that: I was angry and sullen. I never realized anything. I don't dare go a step in such a place. (\textit{They whisper to her again.}) But it's not possible to forgive such things. I don't want to be forgiven so easily. I want to be punished for it all. I won't stir until I've been punished a long, long time. I want to be freed of all that—by punishment. I want to be all new. (\textit{They whisper to her. She puts her feet slowly on the ground.}) But no one else could be punished for me. I'm willing to face it all myself. I don't ask anyone to be punished for me. (\textit{They whisper to her again. She sits long and brokenly looking at her shoes, thinking it over.}) It wasn't fair. I'd have been willing to suffer for it myself—if I could have endured such a mountain. (\textit{She smiles.}) Oh, I'm ashamed! I'm just a stupid and you know it. I'm just another American.—But then what wonderful things must be beginning now. You really want me? You really want me? (\textit{They start leading her down the aisle of the car.})
}

\pair{
    \textbf{哈瑞特}：把这辆车的人都带上吧。这里面有很多好人。我们能一起去吗？你找不到比菲利普更好的人了。求求你们，带我们一起走。
}{
    HARRIET: Let's take the whole train. There are some lovely faces on this train. Can't we all come? You'll never find anyone better than Philip. Please, please, let's all go.
    \newline
    \newline
}

\pair{
    \textit{他们走到楼梯旁，大天使们搭着手，让哈瑞特能够靠在他们身上，慢慢走上楼梯。她听起来像在歌唱，又像在含糊不清说着什么。}
}{
    \textit{They reach the steps. The Archangels interlock their arms as a support for her as she leans heavily on them, taking the steps slowly. Her words are half singing and half babbling.}
    \newline
    \newline
}

\pair{
    \textbf{哈瑞特}：多么高，多么远啊，我的心脏很脆弱，我不应该爬楼梯的。“我不需要看遥远的风景，多走一步就够了。”这里好像瑞士。我无法控制我自己不停地说话——让我停一下脚步，我想要告别。（\textit{在大天使的臂弯里转过身}）就一分钟，我想靠在你的肩上哭一会儿。（\textit{将前额埋在加百列的肩上，转而笑得温柔而绵长}）再见了菲利普——我真希望他没有娶我，但他娶了，他就像你们一样信任我。再见了，伊利诺伊州奥克斯伯里，里奇伍德大道1312号，希望我会永远记得那里的楼梯、房门和墙纸。再见了，弗布什街和维尔利街街角的艾莫森语法学校。再见了，我的英语老师沃克小姐和克莱默小姐，我的生物老师马修森小姐。再见了，梅耶森大道第六街转角的公理会教堂，麦克莱迪博士、麦克莱迪夫人、茱莉亚。再见了，爸爸妈妈……
}{
    HARRIET: But look at how tremendously high and far it is. I've a weak heart. I'm not supposed to climb stairs. “I do not ask to see the distant scene: One step enough for me.” It's like Switzerland. My tongue keeps saying things. I can't control it. —Do let me stop a minute: I want to say good-bye. (\textit{She turns in their arms.}) Just a minute, I want to cry on your shoulder. (\textit{She leans her forehead against Gabriel's shoulder and laughs long and softly.}) Good-bye, Philip. —I begged him not to marry me, but he would. He believed in me just as you do. —Goodbye, 1312 Ridgewood Avenue, Oaksbury, Illinois. I hope I remember all its steps and doors and wallpapers forever. Good-bye, Emerson Grammar School on the corner of Forbush Avenue and Wherry Street. Good-bye, Miss Walker and Miss Cramer who taught me English and Miss Matthewson who taught me biology. Good-bye, First Congregational Church on the corner of Meyer-son Avenue and Sixth Street and Dr. McReady and Mrs. McReady and Julia. Good-bye, Papa and Mama...
    \newline
}

\pair{
    \textit{她转过身。}
}{
    \textit{She turns.}
    \newline
}

\pair{
    \textbf{哈瑞特}：我不想再说再见了。我从没这样说过话。我长得不好看，也不太敢和别人说话。直到菲利普出现。我明白了，我明白了，我现在全都明白了。
}{
    HARRIET: Now I'm tired of saying good-bye. —I never used to talk like this. I was so homely I never used to have the courage to talk. Until Philip came. I see now. I see now. I understand everything now.
    \newline
}

\pair{
    \textit{舞台监督走上前来。}
}{
    \textit{The Stage Manager comes forward.}
    \newline
}

\pair{
    \textbf{舞台监督}：（\textit{对演员们}）好了，好了——现在让我们看看整个世界，整个太阳系。（\textit{全体演员都出现在舞台边上，舞台监督拍拍手。}）整个太阳系。流浪汉呢？月亮去哪儿了？
}{
    THE STAGE MANAGER (\textit{To the actors}): All right. All right. —Now we'll have the whole world together, please. The whole solar system, please. (\textit{The complete cast begins to appear at the edges of the stage. He claps his hands.}) The whole solar system, please. Where's The Tramp? —Where's The Moon?
    \newline
}

\pair{
    \textit{他脚踏了两下地面，像是乐团的指挥要让大家集中注意力，随即举起了他的手。人类低声诉说着他们脑中的想法，时间们在发表演说，星球们在合唱或者低吟。哈瑞特的声音最终盖过了所有人，她说。}
}{
    \textit{He gives two raps on the floor, like the conductor of an orchestra attracting the attention of his forces, and slowly lifts his hand. The human beings murmur their thoughts; The Hours discourse; The Planets chant or hum. Harriet’s voice finally rises above them all, saying:}
    \newline
}

\pair{
    \textbf{哈瑞特}：“向来未曾如此虚心求主，导我前行，从前骄痴无忌，旧事乞莫重提。”
}{
    HARRIET: “I was not ever thus, nor asked that Thou Shouldst lead me on, and spite of fears, Pride ruled my will: Remember not past years.”
    \newline
}

\pair{
    \textit{舞台监督摆摆手让他们下场。}
}{
    \textit{The Stage Manager waves them away.}
    \newline
}

\pair{
    \textbf{舞台监督}：很好。现在请大家离开舞台。我们已经到了芝加哥南部的恩格伍德站了，看见学校的塔楼了吗？比其他地方的都壮观。
}{
    THE STAGE MANAGER: Very good. Now clear the stage, please. Now we're at Englewood Station, South Chicago. See the university's towers over there! The best of them all.
}

\pair{
    \textbf{一号下铺}（\textit{少女音}）：乘务员，你答应过六点一刻叫我起来的。
}{
    LOWER ONE (\textit{The Maiden Lady}): Porter, you promised to wake me up at quarter of six.
}

\pair{
    \textbf{乘务员}：抱歉女士，昨晚发生了一些糟糕的事情，有一位女士病得很厉害。
}{
    THE PORTER: Sorry, ma'am, but it's been an awful night on this car. A lady's been terrible sick.
}

\pair{
    \textbf{一号下铺}：噢，她好点了吗？
}{
    LOWER ONE: Oh! Is she better?
}

\pair{
    \textbf{乘务员}：没呢，一点也不好。
}{
    THE PORTER: No'm. She ain’t one jot better.
}

\pair{
    \textbf{五号下铺}：小伙子，把你的脚从我脸前面拿开。
}{
    LOWER FIVE: Young man, take your foot out of my face.
}

\pair{
    \textbf{舞台监督}（\textit{如前，扮演五号上铺}）：抱歉啊，女士，我滑了一下——
}{
    THE STAGE MANAGER (\textit{Again substituting for Upper Five}): Sorry, lady, I slipped—
}

\pair{
    \textbf{五号下铺}（\textit{发牢骚，但没有很不愉快}）：要我说，这趟旅途就是侵犯他人行为的连续剧。
}{
    LOWER FIVE (\textit{Grumbling not unamiably}): I declare, this trip's been one long series of insults.
}

\pair{
    \textbf{舞台监督}：就一下下，女士，我马上就离开您的视线。
}{
    THE STAGE MANAGER: Just one minute, ma'am, and I'll be down and out of your way.
}

\pair{
    \textbf{五号下铺}：都没有人帮你缝一缝你的袜子吗？你就这样出门不觉得丢人吗？
}{
    LOWER FIVE: Haven't you got anybody to darn your socks for you? You ought to be ashamed to go about that way.
}

\pair{
    \textbf{舞台监督}：抱歉了，女士。
}{
    THE STAGE MANAGER: Sorry, lady.
}

\pair{
    \textbf{五号下铺}：你这样实在是不适合结婚。这才是你最大的问题。
}{
    LOWER FIVE: You're too stuck up to get married. That's the trouble with you.
}

\pair{
    \textbf{九号下铺}：比尔！比尔！
}{
    LOWER NINE: Bill! Bill!
}

\pair{
    \textbf{七号下铺}：啥？你要啥？
}{
    LOWER SEVEN: Yea? Wha'd'ya want?
}

\pair{
    \textbf{九号下铺}：比尔，一般这种列车，要给乘务员多少小费呀？我都好久没回国了……
}{
    LOWER NINE: Bill, how much d'ya give the porter on a train like this? I've been outta the country so long...
}

\pair{
    \textbf{七号下铺}：妈呀，弗雷德，我也不知道啊。
}{
    LOWER SEVEN: Hell, Fred, I don't know myself.
}

\pair{
    \textbf{乘务员}：芝加哥！芝加哥到了！请大家下车！这列列车已经抵达终点站。
}{
    THE PORTER: CHICAGO, CHICAGO. All out. This train don't go no further.
    \newline
}

\pair{
    \textit{乘客们拥挤着下了车。一群带着拖布水桶的老太太上场，准备打扫车厢。}
}{
    \textit{The passengers jostle their way out and an army of old women with mops and pails enter and prepare to clean up the car.}
    \newline
    \newline
}

\pair{
    ——剧终——
}{
    END OF PLAY
}

\end{paracol}

\part{相关资料}

\chapter{剧本介绍}

这部独幕喜剧以1930年12月从纽约开往芝加哥的希亚瓦萨号列车车厢为背景，传统的时间被暂停，唯一衡量存在的方式只有生死。《希亚瓦萨号列车》带领我们踏上一段隐喻性的火车之旅，穿越美国大地，一群形形色色的旅行者被包裹在车厢中，穿越时空，感受着各种情感。

\section{概述}

\subsection{剧情概要}

《希亚瓦萨号列车》从未在纽约或其他地方进行过专业的首演。它的首演于1932年3月19日在俄亥俄州黄泉镇的安提阿克学院举行。这出“《我们的小镇》的未来肌理”以舞台监督在地板上用粉笔画线开始，这些线代表1930年12月21日从纽约开往芝加哥的希亚瓦萨号列车上的铺位。舞台监督引导火车乘客（六组角色中的第一组）登上舞台。“演员们抬着椅子上场。每个人都临时安排了自己的铺位，在用粉笔标记的位置放两把椅子‘面对面’。”唯一的其他布景是“一个通向视线内外的阳台、桥梁或跑道”和“两段从那里下降到舞台的楼梯。”第一组包括住在一号包厢的主人公哈瑞特；住在二号包厢的哈瑞特的丈夫菲利普；三号车厢的丘吉尔夫人（疯女人）；七号下铺的工程师比尔；九号下铺的另一位工程师弗雷德；一号下铺的一位少女；三号下铺的一位中年医生；以及五号下铺的“一位身材魁梧、和蔼可亲的五十岁妇女”。乘务员帮助乘客安顿下来并准备睡觉。“当神一样的舞台监督不再扮演上铺乘客，而是直接对观众讲话时，反幻觉背景下的现实主义表演突然中断。”在舞台监督无所不知的帮助下，乘客的内心想法被听到了。五号下铺正在检查她送给家人的圣诞礼物。七号下铺（比尔）正在去见他住在加利福尼亚的爱人莉莲。三号下铺（医生）”面无表情地读着一本医学杂志，“不时打断他一句疑问句‘所以呢？’”。一号下铺抱怨她的热水袋坏了，睡不着觉。九号下铺对他与苏联人做生意毫无结果感到恼火。舞台监督给了乘务员一个发言的机会；他感到尴尬，但还是分享了他对芝加哥的家和人寿保险的看法。

与此同时，在隔间里，菲利普照顾着哈瑞特，她自晚饭后就一直感觉不舒服。在三号包厢里，这位精神病女人不愿被送进精神病院，但她的护理人员说那是一个“可爱的地方”，到处都是花园和玫瑰。哈瑞特叫来门房，告诉他自己病了，需要医生。医生赶到后，宣布哈瑞特已经去世，并去告诉她的丈夫这个消息。

舞台监督走上前去，将观众从微观的车厢引向宏观层面。“现在从地理、气象、天文和神学角度来考虑“车厢的”位置。”格罗弗角、俄亥俄、田野、流浪汉、帕克斯堡、工人和这位工人被请上舞台，讲述相应地理区域、他们的生活或旅程的统计数据。他们每个人还背诵了一位著名作家、诗人或歌曲的诗句。工人是一个鬼魂，是德国人，“在建造高架桥时被杀，希亚瓦萨号的车厢现在正经过这座高架桥。”他宣读了他的“座右铭”，类似于葛底斯堡演说，然后退场。“这位工人”宣布火车信号工作正常，并祝愿旅客一切顺利。机械师走上前来预报天气。

下一组由女性扮演的演员代表着时间，正如舞台监督所解释的那样：“分钟是闲聊者；小时是哲学家；年份是神学家。小时是哲学家，除了十二点钟，他也是一位神学家。” 十点钟是柏拉图；十一点钟是爱比克泰德；十二点钟是圣奥古斯丁。

时间消逝，星辰降临。它们没有台词，而是被指挥发出声音。舞台监督现在召唤所有人上台，将微观世界和宏观世界融合在一起。“来吧，各位。这是地球的声音。”疯女人请他做点什么，但他让所有人安静下来，大天使们从阳台上下来，这是神学层面。加百列和米迦勒从乘客身边走过。疯女人看到他们，请求大天使带她走；然而，他们让她明白，她必须等待。哈瑞特不愿跟随大天使们。他们低声对她说的话可以通过对话推断出来。“我很惭愧跟你们一起去……但这种事不可能被原谅……但没有人能代替我受到惩罚。”当他们开始带她走的时候，她突然顿悟：“我以前从来不会这样说话。我太丑了，以前从来没有勇气说话。直到菲利普来了。现在我明白了。现在我明白了。现在我明白了一切。”

最后，舞台监督召集全体演员。“他像管弦乐队指挥吸引全体人员的注意力一样，在地板上敲了两下，然后慢慢抬起手。人类低声吟唱着他们的想法；时间在交谈；行星在吟唱或哼唱：哈瑞特的声音终于盖过了一切。 ”她背诵了《引领，仁慈的光》中的一段。哈瑞特和宏观演员们被挥手告别，戏剧回到了微观层面。抵达芝加哥后，乘客们推推搡搡，仿佛什么奇怪的事情都没有发生过。当他们下火车时，“一群拿着拖把和水桶的老妇人进来，准备清理车厢。帷幕落下。”

\subsection{批判性分析}

在讨论《希亚瓦萨号列车》时，经常会提到路易吉·皮兰德娄的《六个寻找作者的角色》。怀尔德的舞台监督与皮兰德娄文本中的导演相呼应，因为他们扮演着演员和导演的角色，并且打破了戏剧的现实；但他们并不完全相同。卡斯特罗诺沃指出，“怀尔德软化了监督的角色，使他不像皮兰德娄笔下脾气暴躁的导演那样暴躁，也不那么抵触眼前的戏剧……（他）成为第一批让观众感觉到戏剧的现实和他们生活的现实是交织在一起的真正的美国人/演员之一。”舞台监督对舞台和角色的掌控表明他是一个神一样的人物。“《希亚瓦萨号列车》的舞台监督更接近于圣经第一本书《创世纪》中上帝的创造者形象。”他控制着谁上台、发生什么，他重复的“好吧”反映了上帝在圣经中所说的“一切都好”。如果舞台监督代表上帝，那么剧中人物则代表了各行各业：一对夫妇、一位祖母、一位商人、几位工程师、一位医生和几位护士、一个疯女人和一个搬运工。他们坐在火车上，这不仅是对人生旅程的隐喻，也是“时间呼啸着驶向目的地”的隐喻。火车在轨道上驶向固定的目的地象征着命运，这进一步强调了为什么哈瑞特必须离开，以及为什么疯女人不能被大天使带走，尽管她愿意并渴望离开。这是哈瑞特的时间，也只是她的时间。

汉松·丹（Hansong Dan）表示，剧中哈瑞特的死亡是怀尔德“在寓言叙事中酝酿他最喜欢的主题之一——死亡的道德和审美价值”。大天使和哈瑞特之间的对话推断出了生命的意义，“这是怀尔德希望传达给观众的寓言信息”。至于火车上的其他乘客，他们仍然是盲人，因为“只有疯子和死人才能真正看到世界”。这些人物不仅代表了各行各业，而且突出了唯物主义人生观的问题，因为他们的独白围绕着他们无法带走的东西。医生被他的杂志分散了注意力；五号下铺只专注于她的礼物；比尔不是被爱情蒙蔽了双眼，而是被爱情的“政治”蒙蔽了双眼；一号下铺忙于思考疾病；弗雷德被金钱蒙蔽了双眼。哈瑞特在说完一系列的再见后，终于“明白了一切”，从而将自己从所有人和所有事中解放出来。保罗·利夫顿指出，哈瑞特的瞬间蜕变是“怀尔德所有戏剧中描绘的最完整、最‘表现主义’的魔杖（转变）……她死后从一个胆怯谦逊的病人……变成了一个冲动、外向、健谈的灵魂。”怀尔德的哲学背景，特别是柏拉图的“至福愿景”，在《希亚瓦萨号列车》中得到了充分体现。阿德里安娜·哈克·丹尼尔斯认为，这是“怀尔德的戏剧似乎在努力追求的一个主题”。哈瑞特能够看到微观和宏观世界中难以察觉的美，并更接近上帝。《我们的小镇》中艾米莉的顿悟与哈瑞特的顿悟如出一辙，且基于相同的前提：她意识到自己一生都对周围的人和对话视而不见。怀尔德通过哲学打开了他笔下人物和观众的眼界。《希亚瓦萨号列车》“比其他作品更能预示怀尔德未来的戏剧技巧和人生哲学。”然而，这种哲学是观众必须理解的；正如丹所说：“怀尔德拒绝以哈瑞特充满说教的诗歌朗诵和观众的哭泣来结束这部戏。”相反，舞台监督指示清理舞台，清洁人员进来。“因此，怀尔德通过将舞台恢复到模仿的状态，摧毁了抚慰人心的宗教幻觉……或许是为了让观众摆脱他们希望从剧院中摆脱的任何一种知识分子的自满情绪。”尽管怀尔德可能无法给出生命意义的答案，但该剧却开启了对生命的讨论，首先探讨的是日常生活中有多少东西是看不见的，这是一个令人不安的真相。

\section{进一步讨论}

《希亚瓦萨号列车》于1931年11月5日首次出版，收录于名为《漫长的圣诞晚餐及其他独幕剧》的合集。1932年3月19日，该剧在俄亥俄州黄泉镇的安提阿克学院首次获得演出许可。二战结束后，该剧及合集内的其他剧目主要由业余剧团使用原著或其他打字稿进行演出。1952年，该剧的表演版在伦敦出版。1993年，该剧在广场圆形剧场（Circle in the Square Theater）重新上演，2001年1月，该剧再次在巴尔的摩中心舞台（Center Stage）上演。

与怀尔德的许多戏剧一样，《希亚瓦萨号列车》中的情节由一位舞台监督精心编排，他既是叙述者，又是导演，偶尔还会客串演员。他负责搭建布景、指导动作，并引导观众对戏剧世界展开的视角。首先，他解释道，布景是由阳台、楼梯和脚灯组成的，是一辆从纽约开往芝加哥的豪华车厢；观众要想象这辆车有一条过道和九个包厢，每个包厢都有上下铺。他把演员们叫到舞台上，指挥他们摆放椅子，以模拟包厢内部的布局。

剧中人物鲜有姓名，人们仅凭其铺位位置而知，例如下铺一号，我们后来得知她是一位“少女”，下铺三号，我们发现她是一位医生。只有丘吉尔夫人、在几个关键时刻占据主导地位的疯女人，以及一对年轻夫妇菲利普和哈瑞特，有明确的名字。剧中核心人物哈瑞特在剧中早早去世，但令人动容的是，她以幽灵的身份再次出现，不再是旅程的一部分；丘吉尔夫人，被反复提及为疯女人，由两名侍从陪同，她也经常向侍从抱怨。大多数角色都睡在下铺。舞台监督则假唱少数几位上铺乘客的声音。

戏剧伊始，剧中人物在火车上安顿下来，谈论着旅途中的各种实际问题，比如热水袋漏水或需要服用阿司匹林。随后，在舞台监督的指示下（乘客发言，观众聆听），他们吐露了自己的想法、希望和担忧。五号下铺担心她要带给家人的礼物；七号下铺担心的是他即将在旧金山——最终目的地——遇到的女人；九号下铺担心的是金钱。所有人都关心着人生的日常琐事。就连搬运工也被鼓励吐露自己的想法，但他却觉得自己不配；对他来说，他对房子和人寿保险的担忧似乎太过简单，不值一提。他不愿开口，没有意识到他的担忧与那些需要他服务和关注的乘客是一样的。舞台监督认为搬运工的不安是害羞的表现。

唯一稍有发言的乘客是疯女人，以及年轻夫妇菲利普和他的妻子哈瑞特。哈瑞特在旅程开始后不久就去世了。这三个人都经历了痛苦而艰难的对话（或者像疯女人那样，是独白。她很难与他人交谈）。

哈瑞特一死，旅程就变了。舞台监督拓展了旅程的背景，仿佛哈瑞特的死扩大了戏剧的语境和重要性。“好了，”他说，“关于车厢内部的描述就到此为止。现在就从地理、气象、天文和神学的角度来探讨它的位置。”首先，他从“地理”的角度来思考旅程。乘客们离开他们的卧铺，成为“俄亥俄州格罗弗角”（《 我们的小镇 》中即将重现的一个村庄）、“俄亥俄州帕克斯堡”，甚至“一片田野”的居民，所有这些都是火车途经的乡村的组成部分。这些地区，就像普尔曼车厢的乘客一样，表达着他们的想法和担忧，而其他角色，例如流浪汉和工人（一位德国铁路工人的鬼魂），讲述的艰辛远比火车乘客的艰辛要凄凉。他们都引用了诗歌，并清晰地将这些诗句归于作者。这位“工人”是一名铁路值班员，他担心火车会晚点，想让火车加速行驶。“信号对这列火车来说没问题，”他说。世界仍在运转。

在确立了火车的尘世背景后，舞台监督再次拓展了戏剧的范畴，通过向机械师询问当地天气，从“气象学”的角度思考世界。接下来，舞台监督略微扩展了他既定的考量，从“时间”的角度呈现了旅程。他以三个“美丽的姑娘”——十点、十一点和十二点——的形式引入时间，并评论道：“现在是十二点。分钟是闲聊，小时是哲学家，年份是神学家。”他补充道，十二点也是神学家。每个人都通过引用哲学家（例如爱比克泰德和圣奥古斯丁）的话语来理解时间。对这三位哲学家来说，时间是通过哲学家的话语体现的人类境况；而“钟点们”则将时间视为用人类语言书写的人类时间。

场景最后一次宏大地转换，舞台监督反思这场“天文式”的旅程，向观众呈现那些无言喧闹的行星，它们的嗡嗡声近乎震耳欲聋。如同管弦乐队指挥，舞台监督指挥着所有角色，奏响一曲低语交响曲，涵盖了旅程的方方面面：乘客的低语，列车途经地点的呼喊，被忽视的幽灵和工人们，充满哲思的“时刻”，以及低语的行星。旅程的宇宙和谐和谐。

但这种宇宙的和谐无法持久。并非所有人都能加入合唱。疯女人用人类的绝望和抱怨打破了“地球之声”的和谐：她感到自己一无是处，她的疯狂与天体美妙的声音格格不入。于是，舞台监督再次转向他的最后一类，即“神学思考”之旅。他召唤神圣的世界，尤其是大天使米迦勒和加百列，引导疯女人和哈瑞特进入另一个世界和生活，在那里，她们，这些旅程中的格格不入者，将拥有归属感。她们将成为一个充满与她们相似之人的地方的居民，与火车上那些活着的、有目标的旅行者的旅程背景格格不入。在人类生活中格格不入者在来世找到了家，在那里他们将“继承地球”。一旦她们离开，充满荣耀与和谐的宇宙乐团将再次演奏，列车上的乘客也将抵达他们短暂的目的地——芝加哥。

《希亚瓦萨号列车》是怀尔德思想的交响曲。该剧探讨了怀尔德长篇作品中探讨的诸多问题。正如怀尔德的传记作者吉尔伯特·A·哈里森所言：“（《希亚瓦萨号列车》）预示了桑顿怀尔德在《我们的小镇》和《我们牙齿的皮肤》中运用的技巧和主题。” 在《我们的小镇》中，怀尔德笔下的俄亥俄州格罗弗斯角变成了新罕布什尔州格罗弗斯角。哈里森指出，哈瑞特动人的告别词使她成为“…… 《我们的小镇 》中艾米丽的先驱”。《我们牙齿的皮肤》重新审视了《希亚瓦萨号列车》的哲学时刻，其中《我们牙齿的皮肤》的舞台监督菲茨帕特里克先生评论道，“……夜晚的每个时刻都代表着一位哲学家或伟大的思想家。比如，十一点钟是亚里士多德，九点钟是斯宾诺莎，就是这样。”菲茨帕特里克的评论与《希亚瓦萨号列车》中的舞台监督的评论如出一辙。最后，火车乘客的对话反映了日常生活的平凡，就像《从卡姆登到特伦顿的快乐旅程》中的对话一样，这也是一次人生旅程。在怀尔德的一些较短的独幕剧中，例如《三人三分钟剧》，这种日常生活与更大的形而上学现实并列，例如在《5:25 的火车失事》中。《希亚瓦萨号列车》与怀尔德的许多戏剧一样，将对话与声音、言说与未言说的巧妙结合，呈现出一种超凡脱俗的气质。这种结合创造了一个超越舞台或戏剧的现实，将戏剧的极限推向了极致。

就像在其他剧作中一样，例如《去特伦顿和卡姆登的快乐旅程》，怀尔德将旅程视为人类实际的和隐喻性的交通工具。哈里森说：“桑顿怀尔德自己也经常在外奔波，他将人类的旅程视为穿越时空的无尽运动。对桑顿怀尔德产生影响的是他30年代初的全国巡回演讲——火车旅行、对寄宿公寓和不起眼的旅馆的瞥见、路上的推销员、女服务员、小镇商人、家庭主妇和他们的孩子。”剧中没有一个角色是多余的。如果他们的生活没有在对话或情节中详细揭示，那么每个人都有机会吐露内心的想法和担忧。即使是风景，例如“田野”，也诉说着这片土地及其居民——地鼠、老鼠和蛇——的重要性和丰富的历史。

怀尔德将这部戏剧视为一张地图。在仔细审视火车及其乘客之后，我们逐渐从狭窄的背景中拉近，进入风景之中。随着我们来到天空、太阳系，最终到达星空本身，全景变得更加开阔。当我们将一些平凡人置于世界秩序的宏大框架下审视时，他们的生活便变得非凡非凡。这或许正是怀尔德的要旨：没有人是无足轻重的，然而在浩瀚的宇宙面前，所有人都显得卑微。用吉尔伯特·哈里森的话来说，“剧作家的意思是：在你之前有数百万人，在你之后还会有数百万人；宇宙浩瀚无垠，时间对骄傲和野心漠不关心，在这短暂的时光里尽你所能，这就是你所能做的一切，上帝知道为什么——或者上帝可能并不知道。”

\chapter{剧中古代人物}

\section{圣奥古斯丁}

"Love and do what you will." If you keep silent, keep silent by love: if you speak, speak by love; if you correct, correct by love; if you pardon, pardon by love; let love be rooted in you, and from the root nothing but good can grow.

“去爱并做你想做的事。”如果你们保持沉默，就用爱保持沉默；如果你们说话，就用爱说话；如果你们纠正，就用爱纠正；如果你们原谅，就用爱原谅；让爱在你们心中扎根，从根部只会生长出善。

\section{爱比克泰德}

人们的困扰不是来自事情的本身，而是来自他们对事情的看法。

我们登上并非我们所选择的舞台，演出并非我们所选择的剧本。

\end{document}